% ===================================================================
%  TABELLA nr. 14 — La via. --- Ils commerziants.
%  Traduction 2026 d'apres texto/io, controlee sur texto/fr, texto/ca
%  et texto/oc. Consigne : voir texto/rm/10-tabella-01.tex.
%
%  UNE BOUTIQUE DE CETTE RUE A ETE TENUE PAR DES ROMANCHES DANS TOUTE
%  L'EUROPE, ET C'EST CELLE DONT LE VOCABULAIRE ROMANCHE EST LE PLUS
%  ETRANGER.
%
%  Les alineas 9 et 10 sont une confiserie. Or les confiseurs des
%  Grisons — les « Zuckerbäcker » de l'Engadine et du Bergell — ont
%  quitte leurs vallees pendant trois siecles pour ouvrir des
%  patisseries a Venise, a Vienne, a Varsovie, a Saint-Petersbourg,
%  a Marseille. C'est le metier que ce pays a exporte.
%
%  ET LE VOCABULAIRE DE CE METIER EST TOUT ENTIER RAPPORTE DU DEHORS :
%  « confisiur », « biscuits », « bombons », « tarta », « paun da
%  spezia ». L'emigrant n'a pas emporte ses mots, il a rapporte ceux
%  des autres.
%
%  C'EST L'INVERSE EXACT DU TABLEAU 9, ET LES DEUX ENSEMBLE DONNENT
%  LA REGLE COMPLETE :
%    quand c'est L'OBJET qui voyage, le mot voyage avec lui — « piz »,
%      « lavina » et « muntanella » sont montes de la vallee vers la
%      plaine avec la chose qu'ils nomment ;
%    quand c'est LE TRAVAILLEUR qui voyage, le mot voyage EN SENS
%      INVERSE — le confiseur part avec son tour de main et revient
%      avec le vocabulaire de sa clientele.
%  Un objet emporte son nom ; un homme rapporte celui des autres.
%
%  Les ramoneurs de l'alinea 20 sont le second cas du meme mouvement.
%  « spazzachamin » est italien, et ce sont les vallees alpines qui
%  ont fourni a l'Europe ses ramoneurs. La aussi l'homme est parti et
%  le mot est rentre etranger.
%
%  Le tableau 13 avait dit qu'un vocabulaire suit sa clientele. Le 14
%  precise ou est la clientele quand c'est l'ouvrier qui bouge : elle
%  est restee la-bas, et c'est elle qui a nomme.
% ===================================================================

%%K t14-apar-1 apar
\VUsaut{4.0mm}
\VUtitre{15.0pt}{60}{TABELLA nr. 14}
\VUsaut{3.0mm}

%%K t14-tit-1 sub
\VUpk{11.4pt}{40}{La via. --- Ils commerziants.}
\VUsaut{3.0mm}

%%K t14-01-1 p
1. --- Mes ami Gion, ch'abitescha en champagna, è vegnì a passentar trais dis en mia
famiglia. \VUgras{Fin} lura n'aveva el mai gì l'occasiun da vesair ina gronda citad;
perquai duvrain nus tut noss dis a spassegiar, e jau al explitgesch quai ch'el
n'enconuscha betg.

%%K t14-02-1 p
2. --- L'emprim essan nus ids a visitar l'\VUgras{observatori} \textsuperscript{(1)}, ch'al
ha fitg interessà. Turnond tras la \VUgras{via} \textsuperscript{(3)} nua ch'è il \VUgras{bogn
tirc} \textsuperscript{(2)}, avain nus scuntrà ina \VUgras{sepultura} \textsuperscript{(4)}. Avant
il \VUgras{cortegi funerar} ha ì il \VUgras{clerus}, che precedeva il \VUgras{char dal
mort}, en il qual era messa la \VUgras{casseta}. Blers amis han accumpagnà la famiglia
\VUgras{en led}.

%%K t14-02-2 p
Suenter il passadi dal cortegi avain nus traversà la via avant la gronda
\VUgras{biereria} \textsuperscript{(5)}, en la quala blers \VUgras{consuments} bevevan
\VUgras{biera} sin la terrassa, protegids da la \VUgras{marchisa} \textsuperscript{(6)} da
vaider.

%%K t14-03-1 p
3. --- Gion era fitg surprais dal \VUgras{stemma} mess tranter il magasin dal
\VUgras{liquorist} \textsuperscript{(7)} e quel dal \VUgras{vendider da musica}
\textsuperscript{(10)}. El m'ha dumandà tge ch'quai signifitgia; jau hai respondì ch'el
mussa il \VUgras{consulat} \textsuperscript{(8)} englais, ed al hai fatg remartgar il
\VUgras{consul} \textsuperscript{(9)}, ch'è sin il balcun.

%%K t14-tit-2 sub
\VUpk{10.2pt}{40}{Guardar en las buttegas.}
\VUsaut{3.0mm}

%%K t14-04-1 p
4. --- Cumprensiblamain essan nus restads avant l'étalage dals \VUgras{instruments da
musica}, \VUgras{harmoniums}, \VUgras{orgas} e \VUgras{pianos}; nus avain guardà ils
cuvertins illustrads da las \VUgras{partituras} e dals \VUgras{tocs da musica} per piano,
ed avain admirà la \VUgras{lira} da lain indorà, che vegn duvrada sco ensaina.

%%K t14-05-1 p
5. --- Ils nivels da pulvra fatgs dals \VUgras{scuvaders} \textsuperscript{(11)} e dal
\VUgras{char da scuvar} \textsuperscript{(12)} ans han sforzads da passar svelt avant las
bellas \VUgras{mobiglias} dal \VUgras{vendider da mobiglias} \textsuperscript{(13)} ed avant
l'étalage da \VUgras{fuorns} e \VUgras{lampas} dal \VUgras{spengler} \textsuperscript{(14)}.

%%K t14-06-1 p
6. --- Dal rest essan nus stads sforzads, en quest lieu, da bandunar il trottuar,
perquai ch'el era ingomber da pachets ch'ins prendeva d'in \VUgras{char da transport}
\textsuperscript{(16)} per als metter en in \VUgras{magasin} \textsuperscript{(15)} ch'ins aveva
gist prendì a fit.

%%K t14-06-2 p
Quai ans ha impedì da vesair ils bels chars dal \VUgras{charrader} \textsuperscript{(17)}, ma
n'ans ha betg impedì da restar per guardar l'\VUgras{impachetader} \textsuperscript{(19)} che
faseva ina chascha per ses vischin, il \VUgras{vendider da velos e d'autos}
\textsuperscript{(18)}. Lura, perquai ch'i era gia in pau tard, essan nus turnads a
chasa.

%%K t14-tit-3 sub
\VUpk{10.2pt}{40}{Spassegiada tras la citad.}
\VUsaut{3.0mm}

%%K t14-07-1 p
7. --- Il di suenter ha mia mamma proponì a Gion ch'ins al fotografeschia, lura m'ha
ella incumbensà da l'accumpagnar tar il \VUgras{fotograf} \textsuperscript{(20)}. Suenter il
gentar essan nus ids en l'atelier da l'artist, nua ch'nus essan stads sforzads da
spetgar bastantamain ditg. Per ans occupar avain nus guardà ils \VUgras{portrets}
\textsuperscript{(22)} che pendan vi dal mir.

%%K t14-07-2 p
Cur ch'i è stà nossa giada, n'ha la lavur betg durà ditg, e gist cur ch'mes ami ha
finì da posar avant l'\VUgras{apparat} \textsuperscript{(21)}, essan nus ids giu.

%%K t14-08-1 p
8. --- Sin l'emprim plaun avain nus vis tras la porta dal \VUgras{coiffeur}
\textsuperscript{(23)}, ch'era averta, ses giuven che tundreva ils chavels d'in client.

%%K t14-08-2 p
Arrivads en la via essan nus passads avant la \VUgras{vendita da tubak}
\textsuperscript{(24)}. Gion è stà uschè divertì da la cotschna \VUgras{lantiarna}
\textsuperscript{(25)} e da la \VUgras{gronda pippa} che vegn duvrada sco \VUgras{ensaina}
\textsuperscript{(26)}, ch'el è ì cunter dus vegls signurs che discurrivan
\VUgras{sfluschond tubak} \textsuperscript{(27)}.

%%K t14-tit-4 sub
\VUpk{10.2pt}{40}{La confisaria.}
\VUsaut{3.0mm}

%%K t14-09-1 p
9. --- Ils \VUgras{bigliets da banca} e las \VUgras{munaidas} da tut ils pajais exponids
en la \VUgras{vitrina} dal \VUgras{cambiader} \textsuperscript{(29)} han tratg per in mument
nossa attenziun, ma perquai ch'i era l'ura da marendar, essan nus entrads tar il
\VUgras{confisiur} \textsuperscript{(31)} senza restar pli ditg.

%%K t14-10-1 p
10. --- Vesend tut las \VUgras{lecheretgas} ch'il magasin cuntegneva, \VUgras{tuortas,
paun da spezia, biscuits} e \VUgras{bombons} da tut las sorts, n'avain nus betg pudì
tscherner facilmain. A la fin ha Gion tschernì las \VUgras{tuortas cun roma}, e jau hai
prendì mo ina pitschna \VUgras{tarta da ceriaschas}, perquai ch'las lecheretgas
\VUgras{fan mal a mes dents} e ch'jau na vi propi betg ir memia savens tar il
\VUgras{dentist} \textsuperscript{(30)}.

%%K t14-tit-5 sub
\VUpk{10.2pt}{40}{Scriver a maschina.}
\VUsaut{3.0mm}

%%K t14-11-1 p
11. --- Per mussar a mes ami ina \VUgras{maschina da scriver}, da la quala el aveva udì
discurrer, ma ch'el n'enconuscheva betg, essan nus entrads en il \VUgras{biro}
\textsuperscript{(32)} da mes \VUgras{cusrin} \textsuperscript{(34)}. Nus al avain chattà
dictond sias brevs a ses \VUgras{secretari} \textsuperscript{(35)}, ch'è \VUgras{stenograf}.
Apaina ch'ina brev era finida, l'ha in \VUgras{dactilograf} \textsuperscript{(33)} copià cun sia
maschina. Na vulend betg interrumper las occupaziuns da mes parent, essan nus restads
tar el mo intgins mumaints.

%%K t14-12-1 p
12. --- Sut il biro da mes cusrin abitescha in grond \VUgras{urer-gioghelier}
\textsuperscript{(37)}, ch'è plinavant aurifiber. Ses splendid magasin cuntegna bleras
\VUgras{uras, pendulas, uras da giaglioffa, chadainettas} e preziusas \VUgras{gioghels},
per exempel \VUgras{culieras da perlas}, \VUgras{bratschadis} d'aur, \VUgras{urengias} ed
\VUgras{anels} ornads cun \VUgras{peidras preziusas} e \VUgras{diamants}. Ina da las
vitrinas è resalvada spezialmain per las \VUgras{ovras d'aur} e cuntegna ina impurtanta
collecziun da \VUgras{posadas} d'argient.

%%K t14-13-1 p
13. --- Virond a l'chantun da la via hai jau remartgà ch'ins aveva midà la
\VUgras{tabla} \textsuperscript{(36)} messa sper il \VUgras{numer} da la chasa. Jau vuleva gist
dir a vusch auta mia remartga, cur ch'ins ha udì ils allegers tuns da \VUgras{la}
musica militara. Nus essan currids encunter al regiment, senza patratgar ch'\VUgras{el}
vegn a passar avant ils magasins dal \VUgras{chapelier} \textsuperscript{(39)} e dal
\VUgras{guantier} \textsuperscript{(40)}, e ch'nus fissan stads precis uschè bain plazzads per
vesair sia defilada restond avant la vitrina da l'\VUgras{opticist}
\textsuperscript{(38)}, plaina da \VUgras{binoclas dal nas}, \VUgras{lorgnons} e
\VUgras{cannocchials}.

%%K t14-tit-6 sub
\VUpk{10.2pt}{40}{Defilada da marsch.}
\VUsaut{3.0mm}

%%K t14-14-1 p
14. --- Avant il \VUgras{regiment} \textsuperscript{(41)} ha ì il \VUgras{chaptambur}
\textsuperscript{(42)}, suandà dals \VUgras{tamburiners} \textsuperscript{(43)}, dals
\VUgras{clariners} \textsuperscript{(44)} e da tut la \VUgras{musica}. Il \VUgras{general}
\textsuperscript{(45)}, ch'era sin la plazza cun ses adjutant, era restà sper il
\VUgras{sgrizzel d'aua} \textsuperscript{(49)} da la \VUgras{funtauna} \textsuperscript{(48)}, per
assistar a la \VUgras{defilada}.

%%K t14-14-2 p
\VUgras{Ils uffiziers dal stab} \textsuperscript{(46)}, il \VUgras{colonel} ed il
\VUgras{lieutenant-colonel} l'han salidà cun la spada. \VUgras{Ils majurs} eran avant lur
\VUgras{bataglions} \textsuperscript{(47)} e mintga \VUgras{capitani} commandava sia
\VUgras{cumpagnia}, entant ch'ils \VUgras{lieutenants}, \VUgras{sutlieutenants} e
\VUgras{sutuffiziers} occupavan lur plaz usità sper lur umens.

%%K t14-15-1 p
15. --- Blers passants s'eran radunads sin \VUgras{la crusch da vias}
\textsuperscript{(50)}; ils commerziants, il \VUgras{paraplievgier} \textsuperscript{(51)}, il
\VUgras{tintader} \textsuperscript{(53)}, l'\VUgras{armarel} \textsuperscript{(54)} èn sortids per
vesair ils \VUgras{schuldads}. Tut ils vehichels, \VUgras{fiakers} \textsuperscript{(52)},
\VUgras{chars a mast} \textsuperscript{(57)} ed era la \VUgras{botta d'arusar}
\textsuperscript{(56)}, èn sa garads lung il trottuar.

%%K t14-tit-7 sub
\VUpk{10.2pt}{40}{Scenas en la via.}
\VUsaut{3.0mm}

%%K t14-15-2 p
In \VUgras{viagiader da commerzi} \textsuperscript{(55)} che purtava en maun sias
\VUgras{chaschas da mustras} ha traversà svelt la via, e las persunas che sesevan sin
l'\VUgras{imperiala} \textsuperscript{(59)} da l'\VUgras{omnibus} \textsuperscript{(58)} sa han
volvidas per giudair il spectacul. In pauper \VUgras{chantadur ambulant}
\textsuperscript{(60)} cun sia \VUgras{orga a maniv} \textsuperscript{(61)} ha stuì interrumper
sia \VUgras{chanzun}, perquai ch'il canaster dals tamburs cuvriva sia vusch.

%%K t14-16-1 p
16. --- Cur ch'il regiment è arrivà avant il \VUgras{commerzi da stoffas}
\textsuperscript{(64)}, han las \VUgras{cusidras} \textsuperscript{(63)} ed ils \VUgras{sarturs}
\textsuperscript{(62)} bandunà lur \VUgras{maschinas da cuser} e lur lavur per guardar tras
la fanestra. \VUgras{Il commerziant} \textsuperscript{(65)}, ch'aveva gist mess novs
\VUgras{custums} \textsuperscript{(66)} vi da ses \VUgras{étalage}, ha stuì far metter dadens
ils \VUgras{draps} \textsuperscript{(67)} e las \VUgras{telas} installadas sin il trottuar,
sper l'\VUgras{avertura da la chanalisaziun} \textsuperscript{(68)}, cun tema ch'la fulla als
rebaltia; schizunt in vegl \VUgras{colportader} \textsuperscript{(69)}, dal qual ina portiera
martgadegiava chaltschas, è stà sforzà d'entrar en in corridor per finir sia vendita.

%%K t14-16-2 p
Nus avain accumpagnà il regiment fin a la chaserna \VUgras{e,} fitg cuntents da nossa
giornada, essan nus turnads a l'ura da la tschaina.

%%K t14-tit-8 sub
\VUpk{10.2pt}{40}{Differents commerzis e professiuns.}
\VUsaut{3.0mm}

%%K t14-17-1 p
17. --- Il di suenter ha mes bab proponì a nus da visitar la stamparia da la gasetta
locala. Nus avain acceptà cun entusiasm.

%%K t14-17-2 p
Il \VUgras{stampader} è era \VUgras{libriar} \textsuperscript{(70)} \VUgras{(= vendider da
cudeschs)}, \VUgras{editur}, \VUgras{vendider da palpiri} e \VUgras{liader da cudeschs}.
El ha installà sin l'emprim plaun la \VUgras{redacziun} \textsuperscript{(71)} da la gasetta,
ed era l'\VUgras{atelier da grava}, en il qual lavuran ils \VUgras{litografs}
\textsuperscript{(72)} ed ils \VUgras{gravaders en rom} \textsuperscript{(73)}, a la fin
l'\VUgras{atelier da cumposiziun}, en il qual èn ils \VUgras{lavurants da stampa}
\textsuperscript{(74)}. La \VUgras{stamparia} \textsuperscript{(75)} propriamain ditga, las
\VUgras{pressas da stampa} e las differentas maschinas èn sin il plaunterren.

%%K t14-tit-9 sub
\VUpk{10.2pt}{40}{Gion fa fitg bleras dumondas.}
\VUsaut{3.0mm}

%%K t14-18-1 p
18. --- Sortind da la stamparia è Gion, fitg surprais, restà avant ina pitschna
chascha da vaider messa sin ina colonna, encunter la \VUgras{merceria}
\textsuperscript{(77)}, ed ha dumandà per tge ch'ella serva. Mes bab al ha explitgà
ch'quai è in \VUgras{avisader d'incendi} \textsuperscript{(76)}. Dal rest era tut, en la
citad, per mes ami in object da surpraisa, da la simpla \VUgras{funtauna da crappa}
\textsuperscript{(78)} ed il ledi \VUgras{triciclo da purtar} \textsuperscript{(80)} manà d'in
\VUgras{grum} \textsuperscript{(79)} cun livrea, fin al \VUgras{char da posta}
\textsuperscript{(81)}.

%%K t14-19-1 p
19. --- Entant che mes bab ans ha bandunads in mument per discurrer cun il
\VUgras{percepter} \textsuperscript{(83)}, ch'era restà per discurrer cun in \VUgras{advocat}
\textsuperscript{(82)} ed in \VUgras{notar} \textsuperscript{(84)}, avain nus ans divertì
suandond cun ils egls la circulaziun sin la via. Las persunas las pli differentas èn
passadas avant nus: in \VUgras{giuven dal confisiur} \textsuperscript{(85)}, che purtava en
in chanaster ina splendida \VUgras{pastetta}, è vegnì suenter in \VUgras{vendider da
vestgadiras} \textsuperscript{(86)}; in \VUgras{emprendider da lantiarnas}
\textsuperscript{(87)} discurriva cun in \VUgras{impiegà dal gas} \textsuperscript{(88)}; ina
\VUgras{sora da convent} \textsuperscript{(89)} che tegneva en maun ina orfna turnava en ses
convent.

%%K t14-tit-10 sub
\VUpk{10.2pt}{40}{Durant noss return a chasa.}
\VUsaut{3.0mm}

%%K t14-20-1 p
20. --- En il mument che mes bab ans ha puspè cuntanschids, ha Gion ris ferm, vesend
la fatscha sporca e la curiusa vestgadira da dus \VUgras{spazzachamins}
\textsuperscript{(91)} dal tut plains fulin.

%%K t14-21-1 p
21. --- Jau vuleva cumprar da la \VUgras{vendidra da flurs} \textsuperscript{(92)} ina
planta per mia mamma, ma mes bab m'ha fatg remartgar ch'ella fiss memia pesanta da
purtar e ch'i fiss meglier da cumprar \VUgras{oranschas}. Nus ans essan avischinads a la
\VUgras{vendidra} \textsuperscript{(94)} e, cur ch'l'\VUgras{educatura} \textsuperscript{(93)},
ch'era per tscherner per ses scolar, ha finì sia cumpra, avain nus laschà servir a
nus.

%%K t14-22-1 p
22. --- Turnond a chasa avain nus scuntrà pliras enconuschientas: ina veglia amia da
mia famiglia, che pusava il bratsch sin sia \VUgras{dama da cumpagnia}
\textsuperscript{(95)}, l'\VUgras{inschigner} \textsuperscript{(96)} da las punts e vias, ed
ina da mias cusrinas cun sia \VUgras{guardgiuffants} \textsuperscript{(98)}.

%%K t14-22-2 p
Lura avain nus scuntrà la \VUgras{chapeliera} \textsuperscript{(97)} da mia mamma e dus
\VUgras{muncs} \textsuperscript{(99)} esters a la citad, ch'èn vegnids tar nus per dumandar
mes bab davart la via a la staziun.

\VUsaut{6.0mm}
\VUfilet{22mm}
